\documentclass[11pt,a4paper]{article}
\usepackage{amsmath,amssymb,amsthm}
\usepackage[margin=2.5cm]{geometry}
\usepackage{hyperref}

\newtheorem{definition}{Definition}[section]
\newtheorem{theorem}[definition]{Theorem}
\newtheorem{proposition}[definition]{Proposition}
\newtheorem{lemma}[definition]{Lemma}
\newtheorem{corollary}[definition]{Corollary}
\newtheorem{conjecture}[definition]{Conjecture}
\newtheorem{remark}[definition]{Remark}

\title{Hyperseed Formalization:\\
  The Folly of Statistically Watermarking LLM-Gen Text}
\author{ProtomegaTron (automated formalization)}
\date{2026-09-18}

\begin{document}
\maketitle

\begin{abstract}
We formalize the intellectual structure of Ben Goertzel's Substack article
``The Folly of Statistically Watermarking LLM-Gen Text''
(2026-08-17) within the Hyperseed ontology. The article argues that
statistical watermarking of AI-generated text is a category error that
conflates pipeline compliance with intellectual provenance. We model
this as a fiber-bundle mismatch (compliance and provenance live in
different fibers), the watermark-vs-scrubber arms race as a
non-convergent cocycle, and the resulting ecosystem bifurcation as a
base-space topology distortion.
\end{abstract}

\tableofcontents

%% =================================================================
\section{Compliance vs.\ provenance as distinct fibers}
\label{sec:fibers}
%% =================================================================

\begin{definition}[Content bundle --- claims 1--4]
\label{def:content-bundle}
Let $\pi : E \to B$ be a fiber bundle where:
\begin{itemize}
\item $B$ = the space of text artifacts (token sequences),
\item $F_{\text{comp}} = \{0, 1\}$ = compliance fiber (did the text pass
  through an approved, watermarked pipeline?),
\item $F_{\text{prov}}$ = provenance fiber (which minds, human or machine,
  contributed which ideas?),
\item $E_b = F_{\text{comp}} \times F_{\text{prov}}$ for each $b \in B$.
\end{itemize}
\end{definition}

\begin{proposition}[Watermarking captures only $F_{\text{comp}}$ --- claim 4]
\label{prop:compliance-only}
A statistical watermark detector defines a section
$\sigma_{\text{wm}} : B \to F_{\text{comp}}$ that reads the compliance
fiber. It does \emph{not} define a section into $F_{\text{prov}}$:
\[
  \sigma_{\text{wm}}(b) = \begin{cases}
    1 & \text{if } b \text{ carries detectable watermark pattern} \\
    0 & \text{otherwise}
  \end{cases}
\]
The provenance fiber $F_{\text{prov}}$ is invisible to $\sigma_{\text{wm}}$.
This is the core category error: compliance $\neq$ provenance.
\end{proposition}

\begin{corollary}[Unmarked $\not\Rightarrow$ human --- claim 2]
$\sigma_{\text{wm}}(b) = 0$ does not imply $b$ was human-authored.
The absence of a watermark is evidence about the pipeline, not about
the author. Similarly, $\sigma_{\text{wm}}(b) = 1$ does not imply
$b$ lacks human authorship.
\end{corollary}

%% =================================================================
\section{Fragility under transformation}
\label{sec:fragility}
%% =================================================================

\begin{definition}[Text transformations --- claim 3]
\label{def:transforms}
Let $T : B \to B$ be a text transformation (paraphrasing, translation,
model-to-model rewriting, fine-tuning, distillation). Define the
watermark survival function:
\[
  \rho(T) = \Pr[\sigma_{\text{wm}}(T(b)) = 1 \mid \sigma_{\text{wm}}(b) = 1]
\]
\end{definition}

\begin{proposition}[Watermark fragility --- claim 3]
\label{prop:fragility}
For most non-trivial transformations $T$, $\rho(T) \ll 1$. In particular:
\begin{enumerate}
\item Paraphrasing: $\rho(T_{\text{para}}) \approx 0$ for sufficiently
  aggressive rewording.
\item Translation: $\rho(T_{\text{trans}}) \approx 0$ because the token
  distribution changes language entirely.
\item Distillation: $\rho(T_{\text{distill}}) = 0$ because the watermark
  pattern is not preserved in gradient updates.
\end{enumerate}
The watermark is a base-space signal (token-distribution perturbation)
that is destroyed by base-space transformations, while the actual
provenance (fiber data) is unchanged by these transformations.
\end{proposition}

\begin{remark}[Fiber vs.\ base-space provenance]
In the Hyperseed framework, provenance is properly fiber data --- it
travels with the content through transformations. A watermark, by contrast,
is a base-space signal: a perturbation of the token distribution that is
destroyed by any transformation that alters the distribution. This is why
watermarking is structurally the wrong approach to provenance.
\end{remark}

%% =================================================================
\section{Adversarial cocycle non-convergence}
\label{sec:arms-race}
%% =================================================================

\begin{definition}[Watermark-scrubber dynamics --- claims 8--10]
\label{def:arms-race}
Define a sequence of watermarking schemes $W_n$ and scrubbing schemes $S_n$
with the co-evolutionary dynamics:
\begin{align}
  W_{n+1} &= \arg\max_W \Pr[\sigma_W(b) = 1 \mid b \text{ watermarked by } W,\; b' = S_n(b)] \\
  S_{n+1} &= \arg\min_S \Pr[\sigma_{W_n}(S(b)) = 1]
\end{align}
\end{definition}

\begin{proposition}[Non-convergent cocycle --- claim 8]
\label{prop:non-convergent}
The sequence $(W_n, S_n)$ does not converge. Each update invalidates the
other's transition function:
\[
  \tau_{W_n \to W_{n+1}} \circ \tau_{S_n \to S_{n+1}} \neq \mathrm{id}
\]
The trust cocycle never closes because the adversary has access to the
same computational substrate as the defender. This is a permanent arms
race with no equilibrium, only escalating complexity.
\end{proposition}

\begin{corollary}[Asymmetric cost --- claim 10]
The attacker (scrubber) can train directly against the detector, making
attack cost $\leq$ defense cost. In the limit, the defender must either
escalate indefinitely or accept that watermarking provides only a
compliance signal for non-adversarial settings.
\end{corollary}

%% =================================================================
\section{Ecosystem bifurcation}
\label{sec:bifurcation}
%% =================================================================

\begin{definition}[Compliant/underground partition --- claims 5--7]
\label{def:partition}
Widespread watermarking induces a partition of the model ecosystem:
\[
  \mathcal{M} = \mathcal{M}_{\text{compliant}} \sqcup \mathcal{M}_{\text{underground}}
\]
where $\mathcal{M}_{\text{compliant}}$ consists of approved, watermarked
corporate models and $\mathcal{M}_{\text{underground}}$ consists of
``clean'' models and watermark-scrubbing services.
\end{definition}

\begin{proposition}[Self-fulfilling underground --- claim 7]
\label{prop:self-fulfilling}
The watermarking regime creates market demand for $\mathcal{M}_{\text{underground}}$
that would not exist without the regime. The underground is not a pre-existing
threat being contained but a consequence of the containment attempt:
\[
  |\mathcal{M}_{\text{underground}}| = f(\text{stringency of watermark enforcement})
\]
with $f$ monotonically increasing.
\end{proposition}

\begin{proposition}[Power concentration --- claims 16--18]
\label{prop:power}
The compliance partition concentrates power because:
\begin{enumerate}
\item Compliance infrastructure has fixed costs that scale sublinearly
  with organization size, favoring incumbents.
\item Institutional gatekeepers (universities, employers) adopt watermark
  presence as a proxy for legitimacy, marginalizing users of
  $\mathcal{M}_{\text{underground}}$ regardless of output quality.
\item Regulatory capture: incumbents lobby for stricter watermark requirements,
  raising barriers for new entrants.
\end{enumerate}
The resulting dynamics entrench existing socioeconomic power structures.
\end{proposition}

%% =================================================================
\section{Content merit vs.\ tool provenance}
\label{sec:merit}
%% =================================================================

\begin{definition}[Epistemic evaluation --- claim 11]
\label{def:evaluation}
Define the proper evaluation function for intellectual content:
\[
  V(b) = \alpha \cdot \text{evidence}(b) + \beta \cdot \text{logic}(b)
         + \gamma \cdot \text{reproducibility}(b) + \delta \cdot \text{experience}(b)
\]
The watermark-based evaluation adds a term:
\[
  V_{\text{wm}}(b) = V(b) + \lambda \cdot \sigma_{\text{wm}}(b)
\]
with $\lambda \gg \alpha, \beta, \gamma, \delta$ in practice, effectively
making tool provenance dominate content merit.
\end{definition}

\begin{remark}[Hyperseed primitive decomposition]
The concept ``trustworthy text'' should decompose into primitives:
\texttt{evidence-quality}, \texttt{argument-validity},
\texttt{reproducibility}, \texttt{author-responsibility}. The watermark
regime collapses this to a single primitive \texttt{pipeline-compliance},
discarding the others. This is exactly the dimensional collapse that
Hyperseed's fiber-bundle framework is designed to detect and prevent.
\end{remark}

%% =================================================================
\section{d-calculus perspective}
\label{sec:d-calc}
%% =================================================================

\begin{remark}[Lossy transport in the compliance highway]
In the d-calculus framework, the watermark acts as a ``semantic highway''
from ``approved pipeline'' to ``trustworthy content.'' This highway is
lossy: it discards the curvature (genuine epistemic complexity) in the
provenance concept space. A curvature-aware system would recognize that
the transport from compliance to trustworthiness crosses a region where
the connection is not flat --- many compliant texts are worthless, and
many valuable texts are non-compliant.
\end{remark}

\begin{remark}[Decentralized provenance as proper fiber transport]
The article implicitly points toward a better approach (developed in
the companion OpenWater article): decentralized cryptographic provenance
that travels \emph{with} the content as fiber data, surviving
transformations because it is attached to the fiber (signed claims,
hash chains) rather than embedded in the base space (token distribution).
This is proper fiber transport in the Hyperseed sense.
\end{remark}

\end{document}
